\begin{activity}{Archean Heat Production}

\section*{Purpose}

To connect radioactive decay, crustal differentiation, and geological processes to the observed heat production of ancient continental crust, we'll investigate the following question ``Why is Archean heat production so low?''

\section*{Radiogenic decay}
The concentration of a radioactive isotope is given by
\[
C(t) = C_0 e^{-\lambda t},
\]
and the corresponding heat production is
\[
A(t) = \sum_i H_i \lambda_i C_i(t),
\]
where $H_i$ is the heat released per decay and $\lambda_i$ is the decay constant.

\section*{Observation}
Simple decay arguments predict higher heat production in the Archean than today, yet preserved Archean crust shows relatively modest values.


\section*{Tasks}
For each mechanism below, explain qualitatively how it could reduce the heat
production of preserved crustal rocks:
\begin{enumerate}
  \item Differentiation and melt extraction of LILEs (K, Th, U).
  \item Selective erosion of felsic upper crust.
  \item Reworking: remelting and redistribution of existing crust.
  \item Foundering of dense mafic lower crust or arc roots.
  \item Secular evolution of mantle temperature.
\end{enumerate}

\subsection*{Differentiation}
Concept: Heat-producing elements (K, Th, U) behave as incompatible elements and strongly partition into melts.

Let $D \ll 1$ be the bulk partition coefficient of a heat-producing element.
For fractional melting,
\[
C_{\text{solid}} = C_0 (1 - F)^{(1/D - 1)},
\]
where $F$ is the melt fraction.  You could think of this simply as:
\[
C_{\text{crust}} \approx \frac{C_{\text{mantle}}}{D} \, F
\]

% Interpretation

% Small melt fractions efficiently extract LILEs into crust.
% The mantle is depleted, and crustal heat production becomes front-loaded early in Earth history.
% This explains early high HP but does not guarantee preservation.

\subsection*{Selective erosion}
Concept: Heat production is concentrated in felsic upper crust, which is preferentially eroded.

Treat the crust as two reservoirs.  Let $A_u$ be upper crust heat production.  If erosion removes heat-producing material at rate $E$,
\[
\frac{dA_u}{dt} = -E A_u.
\]
Solution:
\[
A_u(t) = A_{u0} e^{-Et}
\]

% Interpretation

% Erosion selectively removes the most heat-producing part of the crust.
% Older crust tends to retain lower average heat production, even if originally enriched.
% This acts independently of radioactive decay.

\subsection*{Reworking}
Reworking here will be defined as reheating---past the closure temperatures of geochronometers---and remelting of existing crust rather than net addition of new crust.

Concept: Old crust is repeatedly melted, mixed, and redistributed during later tectonic events.

Schematic formulation: Model reworking as a mixing process.  Let a fraction $R$ of crust be reprocessed per event.
\[
A_{n+1} = (1 - R) A_n + R A_{\text{new}},
\]

where \(A_{\text{new}}\) reflects heat production typical of newly generated crust.

% Interpretation

% Reworking homogenizes and dilutes extreme HP values.
% High Archean HP signatures are progressively erased.
% Preserved Archean crust is not pristine.

\subsection*{Foundering}
Concept: Dense, mafic, heat-poor material is preferentially removed from the lithosphere.

Let the lower crustal mass is reduced over time.  If mafic lower crust is removed at rate $F_d$,
\[
\frac{dM_{lc}}{dt} = -F_d M_{lc}.
\]

The average crustal heat production becomes:
\[
\bar{A}(t) = \frac{A_u M_u + A_{lc} M_{lc}(t)}{M_u + M_{lc}(t)}.
\]

% Interpretation

% Foundering removes heat-poor material.
% This can make the remaining crust appear felsic, buoyant, and thermally stable.
% It biases what crust survives to the present.

\subsection*{Secular cooling}
Concept: The mantle was hotter early, affecting melt production rather than radiogenic inventory.

Mantle potential temperature declines approximately as:
\[
T_m(t) = T_{m0} e^{-t/\tau},
\]

where \(\tau \sim 1\text{--}2\)~Ga is a thermal timescale.
Melt production gives:
\[
F(t) \propto T_m(t) - T_{\text{solidus}}.
\]

% Interpretation

% Hotter mantle → more melting → more differentiation.
% Early Earth transfers HP to crust efficiently.
% Later Earth is melt-limited, not HP-limited.

\begin{table}[h!]
    \centering
    \begin{tabularx}{\textwidth}{@{}XXXX@{}}
    \toprule
    \textbf{Process} &
    \textbf{Primary Physical Mechanism} &
    \textbf{Effect on Heat-Producing Elements (K, Th, U)} &
    \textbf{Primary Geological Consequence} \\
    \midrule
    Radiogenic decay &
    Exponential decay of radioactive isotopes &
    Monotonic decrease of heat production with time &
    Higher heat production in the past; explains first-order time trend only \\[0.3em]

    Differentiation and melt extraction &
    Incompatible elements partition strongly into melts ($D \ll 1$) &
    Rapid transfer of LILEs from mantle to crust early in Earth history &
    Early enrichment of crust; depletion of mantle \\[0.3em]

    Selective erosion &
    Preferential removal of felsic, heat-rich upper crust &
    Progressive loss of the most heat-producing crustal components &
    Older crust appears cooler than expected from decay alone \\[0.3em]

    Reworking (remelting and redistribution) &
    Repeated partial melting and crustal mixing &
    Dilution and homogenization of extreme heat production signatures &
    Preserved Archean crust is rarely pristine \\[0.3em]

    Foundering of mafic lower crust &
    Removal of dense, heat-poor lower lithosphere &
    Bias toward preservation of felsic, buoyant crust &
    Thermally stable, thick continental lithosphere \\[0.3em]

    Secular mantle thermal evolution &
    Early high mantle temperatures enhance melt production &
    Controls efficiency of LILE extraction rather than inventory &
    Strong early differentiation followed by declining crustal growth \\

    \bottomrule
    \end{tabularx}
\end{table}

\section*{Synthesis Questions}
\begin{enumerate}
  \item Which processes change the \emph{distribution} of heat-producing elements?
  \item Which processes affect what crust is \emph{preserved}?
  \item Why does radioactive decay alone fail to predict observed Archean heat production?
\end{enumerate}

Which processes primarily:
\begin{itemize}
  \item redistribute heat-producing elements?
  \item remove heat-producing material from the crust?
  \item bias what crust survives to the present day?
\end{itemize}


\section*{Reflection}
Why is radioactive decay alone insufficient to explain Archean heat production patterns?

% Archean heat production is low not because the early Earth lacked radioactive elements, but because Earth efficiently rearranged and removed them over time.

\textbox{Concept Takeaway: Heat Production, Curvature, and Preservation}
{
Radiogenic decay alone predicts that heat production should decrease smoothly with time, but observed Archean crust shows much lower heat production than this simple expectation.

The heat production of preserved continental crust reflects the combined effects
of:
\begin{itemize}
  \item \emph{redistribution} of heat-producing elements through differentiation and melting,
  \item \emph{removal} of heat-rich material by erosion and tectonic processes, and
  \item \emph{preservation bias}, where buoyant, felsic, and thermally stable crust is more likely to survive.
\end{itemize}

In steady-state geotherms, these processes matter because crustal heat production controls the \emph{curvature} of the geotherm:
\[
T(z) = T_0 + \frac{q_0}{k}z - \frac{A}{2k}z^2.
\]

Two regions with identical surface heat flow can therefore have very different Moho temperatures, lithospheric thicknesses, and long-term stability depending on how heat-producing elements were redistributed and preserved through time.
}

\end{activity}